\section{Conclusion}

We benchmarked ResNet-50, EfficientNet-B0, and ViT-S/16 on HAM10000 under three pre-declared conditions. Focal loss with class-balanced sampling raised ViT-S/16 sensitivity by 8.2 points and lowered the two CNNs, which reversed the point-estimate ranking at 95\% specificity. The direction was consistent across seeds, but every within-architecture 95\% t interval included zero. The EfficientNet-B0 versus ViT interaction was significant at the seed level; the ResNet-50 versus ViT interaction was not. The AUC ordering was preserved, but the gaps closed under focal loss. Calibration error rose 4.4x for ResNet-50, 3.0x for EfficientNet-B0, and 1.36x for ViT-S/16, so calibration tracked the CNN loss rather than the ViT gain.

The fairness analysis became more specific once the cutoff and the error type were separated. The dark stratum had a higher false positive rate at $L^*$ cutoffs of 55, 58, and 60, but none of the dark-versus-light or dark-versus-medium AUC differences was significant. The false-negative $L^*$ association was significant for all three architectures; the false-positive association was significant only for ResNet-50. Fairness in this setting is an error-type shift with limited precision, not a confirmed dark-skin AUC deficit.

The XAI analysis showed that all three architectures place less attribution inside the lesion on darker images. The mass-error relationship differed by architecture, and the false-negative versus false-positive mass values did not separate consistently. HIBA validation showed a consistent CNN sensitivity loss under focal loss, a seed-dependent ViT gain, and the best external performance under strong augmentation. HIBA contains no Fitzpatrick IV melanomas and no Fitzpatrick V-VI images.

The next steps are a clinical dark-skin evaluation with DDI or PASSION, additional seeds for the six focal and CE cells, and a second attribution method. The clinical dataset is the main requirement for a precise Fitzpatrick V-VI claim. The additional seeds would tighten the interaction, and the second attribution method would test whether the tone gradient is specific to Grad-CAM.
